% !TeX root = ../main.tex
%%% ---------------------------------------------------------------------------
%%% 数据记录与处理
%%%
%%% 这一节演示实验报告的核心技能——**规范地报告带不确定度的数据**：
%%%   · 数值一律走 siunitx：\num{93.21(14)} 会排成 93.21 ± 0.14
%%%     （类文件里设了 separate-uncertainty = true）
%%%   · 有效数字保留到不确定度的第一位非零数字（GB/T 3101、JJF 1059.1）
%%%   · 数字列用 S 列按小数点对齐
%%%   · 原始数据放附录，正文只放处理后的结果
%%%   · 表题在表**上方**，图题在图**下方**
%%% ---------------------------------------------------------------------------
\section{数据记录与处理}
\label{sec:data}

\subsection{原始记录}
\label{sec:data:raw}

15 次运行的逐次测试准确率见\cref{tab:raw}，完整的逐 epoch 曲线数据见
\cref{app:rawdata}。

\begin{table}[htbp]
  \centering
  \caption{各优化器 5 次独立运行的测试准确率}
  \label{tab:raw}
  \begin{tabular}{l *{5}{S[table-format=2.2]}}
    \toprule
    \multirow{2}{*}{优化器} & \multicolumn{5}{c}{测试准确率 / \%} \\
    \cmidrule(lr){2-6}
      & {种子 0} & {种子 1} & {种子 2} & {种子 3} & {种子 4} \\
    \midrule
    SGD-M & 94.83 & 94.61 & 95.02 & 94.74 & 94.90 \\
    Adam  & 93.15 & 93.42 & 93.08 & 93.31 & 93.09 \\
    AdamW & 94.55 & 94.71 & 94.38 & 94.66 & 94.49 \\
    \bottomrule
  \end{tabular}
\end{table}

\subsection{数据处理}
\label{sec:data:process}

以 SGD-M 为例，按\cref{eq:mean-std}：

\begin{equation}
  \label{eq:sgdm-mean}
  \bar{x}\subtext{SGD}
    = \frac{94.83 + 94.61 + 95.02 + 94.74 + 94.90}{5}
    = \qty{94.820}{\percent},
\end{equation}

\begin{equation}
  \label{eq:sgdm-std}
  s\subtext{SGD} = \sqrt{\frac{1}{4}\sum_{i=1}^{5}
    \left(x_i - \bar{x}\subtext{SGD}\right)^2} = \qty{0.152}{\percent},
\end{equation}

由\cref{eq:uncertainty} 得平均值的标准不确定度

\begin{equation}
  \label{eq:sgdm-u}
  u(\bar{x}\subtext{SGD}) = \frac{s\subtext{SGD}}{\sqrt{5}}
    = \qty{0.068}{\percent} \approx \qty{0.07}{\percent}.
\end{equation}

按有效数字规则，结果表述为 $\qty{94.82(7)}{\percent}$。三种优化器的
汇总结果见\cref{tab:summary}。

\begin{table}[htbp]
  \centering
  \caption{各优化器的性能汇总（$n = 5$，误差为平均值的标准不确定度）}
  \label{tab:summary}
  \begin{threeparttable}
    \begin{tabular}{l
                    S[separate-uncertainty, table-format=2.2(2)]
                    S[separate-uncertainty, table-format=2.1(2)]
                    S[separate-uncertainty, table-format=3.0(2)]}
      \toprule
      {优化器} & {测试准确率 / \%} & {达到 90\% 的轮数} & {单轮耗时 / s} \\
      \midrule
      SGD-M & 94.82(7)  & 31.6(9)  & 18(1) \\
      Adam  & 93.21(7)  & 18.2(6)  & 19(1) \\
      AdamW & 94.56(6)  & 19.4(7)  & 19(1) \\
      \bottomrule
    \end{tabular}
    \begin{tablenotes}[flushleft]
      \footnotesize
      \item 括号内为不确定度在末位的值，如 94.82(7) 表示
        $\qty{94.82}{\percent} \pm \qty{0.07}{\percent}$。
    \end{tablenotes}
  \end{threeparttable}
\end{table}

\subsection{收敛曲线}
\label{sec:data:curve}

三种优化器的验证准确率曲线见\cref{fig:curve}，阴影为 5 次运行的
$\pm 1$ 标准差范围。

\begin{figure}[htbp]
  \centering
  \begin{subfigure}[b]{0.47\linewidth}
    \centering
    \placeholder{\linewidth}{4.0cm}{figures/acc-curve.pdf}
    \caption{验证准确率}
    \label{fig:curve-acc}
  \end{subfigure}
  \hfill
  \begin{subfigure}[b]{0.47\linewidth}
    \centering
    \placeholder{\linewidth}{4.0cm}{figures/loss-curve.pdf}
    \caption{训练损失}
    \label{fig:curve-loss}
  \end{subfigure}
  \caption{三种优化器的收敛曲线（阴影为 $\pm 1$ 标准差）}
  \label{fig:curve}
\end{figure}

\begin{caution}
  作图时坐标轴必须标注「量 / 单位」，如 \emph{准确率 / \%}、\emph{时间 / s}
  （GB/T 3101 推荐斜杠形式）。图要导出成矢量格式（PDF），
  matplotlib 用 \lstinline|plt.savefig('x.pdf')| 即可。
\end{caution}
